\chapter{Conclusões e Trabalhos Futuros}\label{cap:conclusao}

Este trabalho desenvolveu um sistema embarcado e autônomo que detecta posturas
sentadas prejudiciais e alerta o usuário por vibração, fundindo na cadeira a
distribuição do peso e a orientação da parte superior das costas medida pelo
colar, sem computador, aplicativo ou nuvem.

Quanto aos objetivos específicos, a variação da referência de postura ereta
entre pessoas e entre colocações do sensor mostrou-se da mesma ordem do desvio
a detectar, o que tornou a calibração individual por botão uma condição de
funcionamento (OE1). O colar e a cadeira foram construídos, alimentados por
bateria e integrados por um enlace sem perda de pacotes nos ensaios (OE2). O
alerta por vibração foi percebido pelos usuários e cessou com a correção da
postura (OE3). Com oito participantes, deixando cada um fora do ajuste dos
parâmetros, o sistema alertou em 58 de 62 segmentos de postura prejudicial,
com alerta em 3\,\% do tempo em postura ereta ou reclinada; com pernas
cruzadas, que a regra não pretende sinalizar, essa fração chegou a 28\,\%
(OE4). A principal contribuição é a fusão: os sensores da cadeira,
inadequados como medidores de força, bastaram como indicadores de contato para
não alertar ao reclinar com o quadril apoiado ou ao levantar-se e para
detectar o escorregamento do quadril.

Os resultados indicam a viabilidade técnica da abordagem: os participantes
das capturas tinham biotipos variados, mais de 150 visitantes da mostra
Eureka usaram o protótipo com detecção consistente, e uma célula de
600\,mAh sustentou o colar pelas cerca de 20\,h da mostra. Os percentuais,
porém, vêm de posturas instruídas, e não se avaliou efeito sobre a saúde.
Permanecem como limitações a projeção isolada da cabeça, as pernas
cruzadas, a fixação do colar e a interação entre a vibração e o sensor.

Como trabalhos futuros, propõem-se validar o sistema com mais participantes,
posturas naturais e uso prolongado, já com a correção da interferência do
motor embarcada; adotar vibração com pausas, menos incômoda; compensar a
deriva lenta da referência e aprimorar a fixação do colar; tratar a projeção
da cabeça com outra posição do sensor ou um segundo sensor; e estudar a
mudança de hábito e os sintomas osteomusculares em uso longitudinal.
